\scanpage{123}
\mspage{- 92 -}
n\textsuperscript{o} 371

% L1–2
\uline{de sous-modules} ; \uline{pour tout} $i$, \uline{soit} $u_i=M_i\to N_i$ \uline{une application} $A$-\uline{linéaire et soit} $u=\bigoplus_{1\leq i\leq n}u_i:M\longrightarrow N$. \uline{Alors le diagramme}
% L3–5
\diagram{\begin{tikzcd}[ampersand replacement=\&,column sep=large,row sep=large]
{}^{\mathrm g}\!\bigotimes_{1\leq i\leq n}(\handwritten{\bigwedge}(M_i)) \arrow[r] \arrow[d,"{{}^{\mathrm g}\!\bigotimes_{1\leq i\leq n}(\handwritten{\bigwedge}(u_i))}"']
 \& \handwritten{\bigwedge}(M)\arrow[d,"{\handwritten{\bigwedge}(u)}"]\\
{}^{\mathrm g}\!\bigotimes_{1\leq i\leq n}(\handwritten{\bigwedge}(N_i))\arrow[r] \& \handwritten{\bigwedge}(N)
\end{tikzcd}}
% L6–7
(\uline{où les flèches horizontales sont les isomorphismes définis dans la prop.7}) \uline{est commutatif}.

% L8–10
Il suffit de voir que les deux applications $ {}^{\mathrm g}\!\bigotimes_i(\handwritten{\bigwedge}(M_i))\longrightarrow \handwritten{\bigwedge}(N)$ définies par le diagramme prennent la même valeur dans chacun des modules $\bigotimes_{1\leq i\leq n}(\overset{m_i}{\handwritten{\bigwedge}}M_i)$, ce qui résulte aussitôt des \unsure{définitions|défintions}.

% L11–14
COROLLAIRE 3.- \uline{Soit} $u:M\to N$ \uline{une application} $A$-\uline{linéaire injective}, \uline{telle que} $u(M)$ \uline{soit facteur direct de} $N$. \uline{Alors l'homomorphisme} $\handwritten{\bigwedge}(u):\handwritten{\bigwedge}(M)\longrightarrow \handwritten{\bigwedge}(N)$ \uline{est injectif et} $\handwritten{\bigwedge}(N)$ \uline{est somme directe de la sous-algèbre graduée} $\handwritten{\bigwedge}(u)(\handwritten{\bigwedge}(M))$ \uline{et d'un idéal bilatère gradué engendré par un supplémentaire de} $u(M)$ \uline{dans} $N$.

% L15–20
Cela résulte du cor.1 appliqué pour $n=2$, car la somme des sous-modules $(\overset{p}{\handwritten{\bigwedge}}\,M_1)\otimes(\overset{q}{\handwritten{\bigwedge}}\,M_2)$ pour tous les couples $(p,q)$ tels que $p\geq 1$ n'est autre que l'idéal bilatère $m_1\otimes(\handwritten{\bigwedge}(M_2))$ de $(\handwritten{\bigwedge}(M_1))\otimes^{\mathrm g}(\handwritten{\bigwedge}(M_2))$, où $m_1$ est la somme directe des $\handwritten{\bigwedge}(M_1)$ pour $p\geq 1$, idéal bilatère de $\handwritten{\bigwedge}(M_1)$. On peut aussi donner une démonstration directe n'utilisant pas la prop.7 et calquée sur celle du \S~4, n\textsuperscript{o}2, prop.4.

% L21–22
COROLLAIRE 4.- \uline{Tout élément} $z$ \uline{de} $\handwritten{\bigwedge}(M)$ \uline{dont la composante homogène scalaire} (i.e. \uline{de degré} 0) \uline{est nulle}, \uline{est nilpotent}.

% L23
Il est clair en effet que tout élément de $\handwritten{\bigwedge}(M)$ appartient à une sous-algèbre
% archive sheet number: 122
